% !TEX root = rechenschaftsbericht.tex
\newcommand{\jahr}{2025}
\newcommand{\todo}[1]{\textcolor{red}{ToDo: #1}}
\input{../vorlage}

% verhindert, dass nicht trennbare Begriffe (deRSE25, ECEASST, URLs) in den Rand laufen
\emergencystretch=3em

\begin{document}
\thispagestyle{empty}

\begin{centering}
\includegraphics[height=3em]{de-RSE-logo-text-colour}\\
\vspace{3em}
\textbf{
 \Large Rechenschaftsbericht des Vorstands\\*[.5em]
 \normalsize Geschäftsjahr: \jahr}\\*[3em]
\end{centering}

\section{Mitglieder und Mitgliedsbeiträge}

Der Verein hatte zu Beginn des Geschäftsjahres \todo{Anzahl} und zu dessen Ende \todo{Anzahl} Mitglieder.

Bei den Mitgliedsbeiträgen wurden zwei Regelungen präzisiert bzw. neu eingeführt: Nachweise für die
Beitragsreduzierung von Studierenden werden auch dann akzeptiert, wenn sie nicht von einer offiziellen
Bildungsstätte ausgestellt wurden. Zudem beschloss die Mitgliederversammlung am 11.~November eine
Änderung der Geschäftsordnung, die für Angehörige von Einrichtungen, die institutionelle Mitglieder des
de-RSE~e.V. sind, einen reduzierten Beitrag von 45~Euro vorsieht. Die Aufnahme institutioneller
Mitglieder selbst steht weiterhin aus, da die dafür notwendige Satzungsänderung noch nicht im
Vereinsregister eingetragen ist.

\section{Vorstand}

Dieser Rechenschaftsbericht wird vom Vorstand des Geschäftsjahres 2025 vorgelegt, welcher sich aus folgenden Personen zusammensetzt:

\begin{itemize}
  \setlength{\itemsep}{0pt plus 1pt}
  \item Jan Linxweiler (Vorsitzender)
  \item Frank Löffler (stellvertretender Vorsitzender)
  \item Jan Philipp Dietrich (Schriftführer)
  \item Bernadette Fritzsch (stellvertretende Schriftführerin)
  \item Stephan Janosch (Schatzmeister)
  \item Michael Meinel (stellvertretender Schatzmeister)
\end{itemize}

Der Vorstand wurde von der Mitgliederversammlung am 11.~November 2025 in unveränderter Besetzung
wiedergewählt. Da die Satzungsänderung zur Verlängerung der Amtsperiode auf zwei Jahre zum Zeitpunkt
der Wahl noch nicht im Vereinsregister eingetragen war, erfolgte die Wahl nach der bisherigen Satzung
und damit für eine Amtsperiode von einem Jahr. Im Zuge der Einführung einer Datenschutzordnung wurde
Michael Meinel im März zum Datenschutzbeauftragten des Vereins bestellt.

\section{Ereignisse im Zeitverlauf}

\begin{itemize}
 \item[] \textbf{1. Quartal}
   \begin{itemize}
     \item Die Klaus Tschira Stiftung (KTS) fördert den Verein im Projekt
       \href{https://www.futursi.de/}{FutuRSI} ("`Konzeption und Anlaufphase eines deutschen
       Forschungssoftware-Instituts"'), das von einem Konsortium aus sechs Institutionen getragen wird.
       Gemeinsam mit der Fachgruppe RSE der GI sind zudem die Projekte "`RSE Master"' und
       "`RSE Award"' erfolgreich. Der Verein
       \href{https://de-rse.org/blog/2025/01/20/kts_foerderung.html}{berichtet darüber im Blog}.
     \item Der gemeinsame Arbeitskreis von GI und de-RSE veröffentlicht die
       \href{https://doi.org/10.18420/2025-gi_de-rse}{\glqq GI- und DE-RSE Muster-Leitlinie zur effizienten
       Entwicklung von Forschungssoftware\grqq}.
     \item Das BMBF lässt eine Studie zu "`Open Research Tools"' erstellen. Bernadette Fritzsch
       vertritt die Interessen des Vereins im Expert*innenworkshop am 22.~Januar.
     \item Vom 25.\ bis 27.~Februar findet am Karlsruher Institut für Technologie die deRSE25-Konferenz
       statt, erstmals gemeinsam mit der Fachtagung Software Engineering (SE25) der GI, mit
       \todo{Anzahl} Teilnehmenden und \todo{Anzahl} Beiträgen. Der Vorstand zeigt sich gemeinsam bei
       der Eröffnung und ist während der Postersession an einem eigenen Stand ansprechbar.
     \item In einer Sondersitzung am 25.~Februar berät der Vorstand über die Weiterentwicklung der
       Veranstaltungsformate des Vereins und legt als grobe Planungsrichtung eine nationale Konferenz
       2026 und eine internationale Konferenz 2027 fest.
     \item Der Verein erhält einen persistenten Identifikator in der Research Organization Registry
       (\href{https://ror.org/007qpef44}{ROR: 007qpef44}) und kann damit in Publikationen eindeutig
       als Affiliation angegeben werden.
     \item Der Vorstand beschließt eine Datenschutzordnung mit Vertraulichkeitsverpflichtung,
       Datenschutzinformation und Datenschutzrichtlinie und bestellt Michael Meinel zum Datenschutzbeauftragten.
     \item Nach Abschluss der Review-Phase bestätigt der Vorstand das
       \href{https://github.com/DE-RSE/2023_paper-RSE-groups}{Papier \glqq Establishing RSE departments in
       German research institutions\grqq} als offizielles Positionspapier des Vereins.
     \item Der Vorstand unterstützt die Zusammenführung der weitgehend inhaltsgleichen Mailinglisten
       des Vereins und der Fachgruppe RSE der GI.
   \end{itemize}\clearpage
 \item[] \textbf{2. Quartal}
   \begin{itemize}
    \item Mit dem Projektstart von FutuRSI im April tritt der Verein erstmals als Arbeitgeber auf.
      Damit verbunden sind der Aufbau der notwendigen Strukturen für Lohnabrechnung, Personal- und
      Datenschutzfragen sowie der quartalsweise Mittelabruf und der jährliche Verwendungsnachweis
      gegenüber der KTS. Für den Projektmitarbeiter wird ein Arbeitsplatz bei der GI in Berlin
      finanziert.
    \item Der Vorstand beauftragt die Steuerkanzlei S\&K mit der Erstellung der Steuererklärungen für
      die Jahre 2025 bis 2027. Die Steuererklärungen für die Jahre 2020 bis 2022 wurden zuvor
      abgeschlossen.
    \item Am 6.\ und 7.~Mai findet der Kickoff-Workshop des Projekts FutuRSI an der SUB Göttingen
      statt, bei dem ein Project Canvas und eine gemeinsame Vision erarbeitet werden.
    \item Am 12.~Mai werden die Ergebnisse der Studie zu Open Research Tools im BMFTR vorgestellt.
      Zu den Empfehlungen zählen die Forcierung der Anerkennung von Forschungssoftware als
      wissenschaftlicher Output und die Würdigung von "`Openness"' als zentrales Element.
    \item Termin und Ort der deRSE26 werden festgelegt: Die Konferenz findet Anfang März 2026 an der
      Universität Stuttgart statt. Ein lokales Organisationsteam nimmt seine Arbeit auf, der Call for
      Committee wird vorbereitet.
    \item Der Arbeitskreis "`Kategorien von Forschungssoftware"' veröffentlicht seine Arbeit als
      \href{https://doi.org/10.1109/MCSE.2025.3555023}{\glqq Multi-Dimensional Research Software
      Categorization\grqq}.
    \item Der Vorstand trifft sich am 20.~Juni zu einer Sitzung in Präsenz in Braunschweig.
   \end{itemize}
 \item[] \textbf{3. Quartal}
   \begin{itemize}
    \item Das Projekt FutuRSI nimmt inhaltlich Fahrt auf: Die Erfassung der RSE-Gruppen in Deutschland
      beginnt, ein Leitfragenkatalog für Stakeholder-Interviews wird erstellt und die Interviews
      werden geführt, ein Kommunikationsplan wird entwickelt. Weitere Projektmitarbeitende bei den
      Partnern nehmen ihre Arbeit auf.
    \item Der Verein kündigt die \href{https://de-rse.org/blog/2025/09/22/deRSE26-call-for-contribution.html}{deRSE26}
      an: 3.\ bis 5.~März 2026 an der Universität Stuttgart, mit einem Call for Contributions und
      einem geplanten Post-Proceedings-Band.
    \item Termin und Ort des RSE Collaborations Workshops 2026 stehen fest und werden an die Community
      kommuniziert: 23.\ bis 25.~September 2026 in Göttingen, mit einem AI4RSE-Workshop am Vortag.
    \item Vom 15.\ bis 16.~September findet das Fachgruppentreffen RSE der GI in Potsdam statt, im
      September zudem das Informatik-Festival der GI, an dem der Verein mit einem Beitrag vertreten
      ist. Ein Treffen zur Rolle von RSEs innerhalb der NFDI findet im Umfeld der CoRDI in Aachen
      statt.
    \item Das Positionspapier zu RSE-Gruppen wird zur Veröffentlichung eingereicht; die überarbeitete
      und finale Version des Papiers
      \href{https://doi.org/10.12688/f1000research.157778.2}{\glqq Foundational Competencies and
      Responsibilities of a Research Software Engineer\grqq} liegt vor.
   \end{itemize}
 \item[] \textbf{4. Quartal}
   \begin{itemize}
    \item Der Arbeitskreis "`Kommunikationsstrategie"' nimmt seine Arbeit auf; seit Oktober erscheint
      regelmäßig ein Newsletter des Vereins.
    \item Am 9.~Oktober beteiligt sich die Community am internationalen RSE Day.
    \item Der Vorstand beschließt am 16.~Oktober einstimmig, aufgrund des aktuellen Finanzrahmens und
      geplanter Ausgaben kein Reisestipendium für den Workshop "`Software Engineering and Research
      Software"' der ICSE~2026 zu vergeben.
    \item Am 11.~November findet die Mitgliederversammlung online statt. Der Vorstand wird für das Jahr
      2024 entlastet, die Änderung der Geschäftsordnung zu reduzierten Mitgliedsbeiträgen angenommen,
      der Vorstand in unveränderter Besetzung wiedergewählt und Thomas Krause als Kassenprüfer sowie
      Alexander Struck als stellvertretender Kassenprüfer gewählt. Die Versammlung beschließt zudem
      den Beitritt des Vereins zur Coalition for Advancing Research Assessment (CoARA).
    \item Am 10.~November findet das Berichtsgespräch mit der KTS zum Projekt FutuRSI statt und
      verläuft erfolgreich.
    \item Am 20.~November richtet der Vorstand den Arbeitskreis "`RSE @ CoARA"' mit Alexander Struck
      als Sprecher einstimmig ein. Im Dezember wird der Antrag auf CoARA-Mitgliedschaft eingereicht;
      damit beginnt eine Frist von einem Jahr, innerhalb derer ein Action Plan vorzulegen ist.
    \item Die Post-Proceedings der deRSE25 erscheinen in der
      \href{https://eceasst.org/index.php/eceasst/issue/view/207}{Zeitschrift ECEASST}.
    \item Am 17.\ und 18.~Dezember veranstaltet der Arbeitskreis "`Teaching RSE"' in Dresden ein
      \href{https://the-teachingrse-project.github.io/teachingRSE-symposium/}{Symposium on Teaching
      Research Software Engineering in Germany}.
    \item Der Vorstand befürwortet, die Mitgliederversammlung 2026 im Rahmen des RSE Collaborations
      Workshops als öffentliche, hybride Veranstaltung abzuhalten.
   \end{itemize}
\end{itemize}
\clearpage
\section{Weitere Ereignisse}

\begin{itemize}
 \item[] \textbf{Blog}
 \begin{itemize}
  \item \href{https://de-rse.org/blog/2025/01/20/kts_foerderung.html}{Klaus Tschira Stiftung fördert de-RSE im Projekt FutuRSI}
  \item \href{https://de-rse.org/blog/2025/02/11/musterleitlinie.html}{GI- und DE-RSE Muster-Leitlinie zur effizienten Entwicklung von Forschungssoftware veröffentlicht}
  \item \href{https://de-rse.org/blog/2025/09/22/deRSE26-call-for-contribution.html}{Call for Contribution -- deRSE26, the RSE Conference 2026 in Germany}
 \end{itemize}
 \item[] \textbf{Vorstandssitzungen}\\
  13 Vorstandssitzungen fanden an den folgenden Tagen statt: 16.1., 20.2., 25.2. (Sondersitzung),
  27.3., 10.4., 15.5., 20.6., 17.7., 21.8., 2.10., 16.10., 20.11., 17.12.
  Alle Sitzungen waren beschlussfähig; in der Sitzung am 2.10. entfiel die Beschlussfähigkeit im
  Verlauf der Sitzung.
  Protokolle sind öffentlich unter \href{https://github.com/DE-RSE/protokolle/}{https://github.com/DE-RSE/protokolle/} einsehbar.
 \item[] \textbf{Mitgliederversammlung}\\
  Die Mitgliederversammlung fand am 11.11. online statt.
 \item[] \textbf{Sprecher*innen-Rat}\\
  In seinem ersten vollständigen Geschäftsjahr tagte der Sprecher*innen-Rat der Arbeitskreise
  viermal: 17.3., 16.6., 22.9., 15.12.
  Im Jahr 2025 waren die folgenden Arbeitskreise aktiv, mehrere davon gemeinsam mit der Fachgruppe RSE
  der GI: Vereinsinfrastruktur, International Council, Positionspapier RSE-Gruppen, NFDI, deRSE25,
  deRSE26, Veranstaltungskoordination, RSE Software Development Guidelines, Kategorien von
  Forschungssoftware, RSE Advocacy Strategy, RSE State of the Nation Report, RSE Research, Teaching
  RSE, AI4RSE, Kommunikationsstrategie sowie ab November RSE~@~CoARA.
 \item[] \textbf{OpenScience/deRSE-Community-Calls}\\
 Mitglieder des Vorstands beteiligen sich an 8 von 12 OpenScience/deRSE-Community-Calls im Jahr 2025.
\end{itemize}

% Mindestinhalt laut https://www.vereinswelt.de/rechenschaftsbericht
%Mitgliederentwicklung: Zu- und Abgang von Mitgliedern, Erläuterungen zu auffälligen Entwicklungen, Ausschlussverfahren
%Durchgeführte Vereinsveranstaltungen
%Teilnahme an Wettbewerben und Ergebnisse
%Beziehungen zum Dachverband und zu anderen Vereinen
%Stand laufender Projekte
%Struktur des Vereins
%Aktivitäten der Organe und Ausschüsse
%Sonstige Ereignisse, die für den Verein wichtig waren
%Finanzbericht
% Empfohlen
%Beziehungen zu Sponsoren und Spendern
%Aktivitäten zur Gewinnung weiterer Sponsoren und Spender
%Ausgang von für den Verein bedeutsamen Gerichtsverfahren
%Hauptamtliche Mitarbeiter, Veränderungen im Personalbestand
%Geplante Projekte und Aktivitäten


\end{document}
